\newpage
\section{Derivadas de estabilidade}

As derivadas foram levantadas para a configuração de CG traseiro
(\texttt{aft.avl}) no ponto de projeto, sem restrição de trimagem: meta de
sustentação $C_L = 0{,}5053$ em $M = 0{,}85$, profundor fixo em zero e
incidência de empenagem de cruzeiro ($i_t = -5{,}90^\circ$), o que resulta
em $\alpha = 4{,}60^\circ$. As derivadas longitudinais e de força lateral
vêm do comando \texttt{st} (eixos de estabilidade), os momentos
látero-direcionais do comando \texttt{sb} (eixos do corpo), e as derivadas
de arrasto em relação a $q$, $i_t$ e $\delta_e$ de diferenças finitas
centradas na mesma condição. Os dados gerais da aeronave estão na
Tabela~\ref{tab:dados_gerais}, com os momentos de inércia calculados pela
função auxiliar da disciplina com a fração de combustível de 44\% que
reproduz o peso do ponto de projeto e 100\% de carga paga.

\begin{table}[H]
\centering
\caption{Informações gerais para análise de estabilidade e controle.}\label{tab:dados_gerais}
\renewcommand{\arraystretch}{1.25}
\begin{tabular}{lcl}
\toprule
Parâmetro & Valor & Fonte \\
\midrule
$S_{\mathrm{ref}}$ [m$^2$] & 368,833 & \texttt{designTool} \\
$c_{\mathrm{ref}}$ [m] & 6,8995 & \texttt{designTool} \\
$b_{\mathrm{ref}}$ [m] & 60,1518 & \texttt{designTool} \\
$m$ [kg] & 229.669,3 & \texttt{designTool} \\
$I_{xx}$ [kg$\cdot$m$^2$] & $1{,}508\times 10^{7}$ & \texttt{designTool} \\
$I_{yy}$ [kg$\cdot$m$^2$] & $4{,}284\times 10^{7}$ & \texttt{designTool} \\
$I_{zz}$ [kg$\cdot$m$^2$] & $5{,}578\times 10^{7}$ & \texttt{designTool} \\
$I_{xz}$ [kg$\cdot$m$^2$] & $1{,}209\times 10^{6}$ & \texttt{designTool} \\
$i_p$ [graus] & 0,0 & \texttt{designTool} \\
$x_p$ [m] & 18,0 & \texttt{designTool} \\
$z_p$ [m] & 3,0 & \texttt{designTool}, sinal invertido para o MVO \\
$T_{\mathrm{max}}$ [N] & 1.026.000 & \texttt{designTool}, 513.000 por motor \\
$V$ [m/s] & 252,1 & ponto de projeto \\
$h$ [m] & 10.668 & ponto de projeto \\
\bottomrule
\end{tabular}
\end{table}

\subsection*{Resultados em $\alpha = 0$ e ajuste da polar}

O $C_{L0}$ e o $C_{M0}$ vêm de uma rodada com $\alpha = 0$, $i_t = 0$ e
profundor nulo. O ajuste quadrático $C_D = C_{D0} + C_{D\alpha}\,\alpha +
C_{D\alpha 2}\,\alpha^2$ foi feito sobre a polar não trimada do CG
traseiro (item 5.b), com $\alpha$ em radianos.

\begin{table}[H]
\centering
\caption{Resultados para $\alpha = 0$, $i_t = 0$ e ajuste quadrático da polar.}\label{tab:cl0_polar}
\renewcommand{\arraystretch}{1.25}
\begin{tabular}{lcl}
\toprule
Parâmetro & Valor & Fonte \\
\midrule
$C_{L0}$ & 0,0670 & AVL \texttt{ft} \\
$C_{M0}$ & $-0{,}1346$ & AVL \texttt{ft} \\
$C_{D0}$ & 0,01954 & ajuste do item 5.b \\
$C_{D\alpha}$ [1/rad] & $-0{,}0424$ & ajuste do item 5.b \\
$C_{D\alpha 2}$ [1/rad$^2$] & 1,5770 & ajuste do item 5.b \\
\bottomrule
\end{tabular}
\end{table}

\subsection*{Tabela completa de derivadas}

A Tabela~\ref{tab:derivadas} apresenta os valores lidos do AVL e os
valores finais para o MVO, já com as conversões do roteiro aplicadas: as
derivadas de controle ($i_t$, $\delta_e$, $\delta_a$ e $\delta_r$) foram
multiplicadas por $180/\pi$, e os sinais de $C_{Y\beta}$, $C_{Yp}$,
$C_{Yr}$, $C_{\ell\delta a}$, $C_{\ell\delta r}$, $C_{n\delta a}$ e
$C_{n\delta r}$ foram invertidos pela diferença de convenções.

\begin{table}[H]
\centering
\caption{Derivadas de estabilidade para o CG traseiro no ponto de projeto.}\label{tab:derivadas}
\renewcommand{\arraystretch}{1.15}
\begin{tabular}{lccl}
\toprule
Parâmetro & Valor AVL & Valor MVO & Fonte e conversão \\
\midrule
$C_{L\alpha}$ [1/rad] & 6,4672 & 6,4672 & \texttt{st} \\
$C_{Lq}$ [1/rad] & 10,8669 & 10,8669 & \texttt{st} \\
$C_{Lit}$ & 0,013865 & 0,7944 & \texttt{st}, $\times 180/\pi$ \\
$C_{L\delta e}$ & 0,007719 & 0,4423 & \texttt{st}, $\times 180/\pi$ \\
$C_{Dq}$ [1/rad] & 0,1535 & 0,1535 & dif. finita \\
$C_{Dit}$ & $-0{,}000170$ & $-0{,}0097$ & dif. finita, $\times 180/\pi$ \\
$C_{D\delta e}$ & $-0{,}000090$ & $-0{,}0052$ & dif. finita, $\times 180/\pi$ \\
$C_{M\alpha}$ [1/rad] & $-2{,}2551$ & $-2{,}2551$ & \texttt{st} \\
$C_{Mq}$ [1/rad] & $-32{,}4502$ & $-32{,}4502$ & \texttt{st} \\
$C_{Mit}$ & $-0{,}052910$ & $-3{,}0315$ & \texttt{st}, $\times 180/\pi$ \\
$C_{M\delta e}$ & $-0{,}030696$ & $-1{,}7588$ & \texttt{st}, $\times 180/\pi$ \\
$C_{Y\beta}$ [1/rad] & $-0{,}4029$ & 0,4029 & \texttt{st}, sinal invertido \\
$C_{Yp}$ [1/rad] & 0,0251 & $-0{,}0251$ & \texttt{st}, sinal invertido \\
$C_{Yr}$ [1/rad] & 0,4386 & $-0{,}4386$ & \texttt{st}, sinal invertido \\
$C_{Y\delta r}$ & $-0{,}004035$ & $-0{,}2312$ & \texttt{st}, $\times 180/\pi$ \\
$C_{\ell\beta}$ [1/rad] & $-0{,}2354$ & $-0{,}2354$ & \texttt{st} \\
$C_{\ell p}$ [1/rad] & $-0{,}5602$ & $-0{,}5602$ & \texttt{sb} \\
$C_{\ell r}$ [1/rad] & 0,1790 & 0,1790 & \texttt{sb} \\
$C_{\ell\delta a}$ & $-0{,}004157$ & 0,2382 & \texttt{sb}, sinal e $\times 180/\pi$ \\
$C_{\ell\delta r}$ & $-0{,}000505$ & 0,0289 & \texttt{sb}, sinal e $\times 180/\pi$ \\
$C_{n\beta}$ [1/rad] & 0,1735 & 0,1735 & \texttt{st} \\
$C_{np}$ [1/rad] & $-0{,}0905$ & $-0{,}0905$ & \texttt{sb} \\
$C_{nr}$ [1/rad] & $-0{,}2037$ & $-0{,}2037$ & \texttt{sb} \\
$C_{n\delta a}$ & $-0{,}000136$ & 0,0078 & \texttt{sb}, sinal e $\times 180/\pi$ \\
$C_{n\delta r}$ & 0,002190 & $-0{,}1255$ & \texttt{sb}, sinal e $\times 180/\pi$ \\
\bottomrule
\end{tabular}
\end{table}

Duas verificações de sanidade sustentam a tabela. A derivada $C_{Lq}$
obtida por diferença finita coincide com a do \texttt{st} em todas as
casas apresentadas, validando o procedimento numérico usado para as
derivadas de arrasto. E a relação $x_{np} = x_{cg} -
c_{\mathrm{ref}}\,C_{M\alpha}/C_{L\alpha}$ reproduz o ponto neutro da
Seção 7 com o $C_{M\alpha}$ e o $C_{L\alpha}$ da tabela.
